\documentclass[11pt,a4paper]{article}
\usepackage[margin=2.1cm]{geometry}
\usepackage{amsmath,amssymb,booktabs,graphicx,hyperref}
\graphicspath{{docs/figures/}}
\hypersetup{colorlinks=true,urlcolor=blue,linkcolor=blue}
\title{Project Briefing: Lorenz PINNs and Residual Direction}
\author{Supervisor-facing status note}
\date{22 November 2026}
\begin{document}
\maketitle

\noindent\textbf{One-minute account.} The project asked whether the
\emph{direction} of a trained physics-informed neural network's differential-
equation residual predicts later error on a chaotic Lorenz trajectory beyond
the residual's magnitude. Its original planned test, using 40 PINNs, did not
support that pointwise hypothesis. The team then improved the PINN solution
itself using eight sequential short networks. Accuracy over two time units
held in 35/40 development runs, then in 30/40 runs on \emph{new untouched
initial conditions}. Repeating the direction test on the new trajectories
again found no reliable gain from one expanding direction. An oracle using
the full tangent propagator did add held-out information at longer horizons.
That is a clue about the mechanism, not yet a practical error indicator.

\section*{Goal and original hypothesis}
The goal was to find out whether residual orientation could provide a useful
\emph{post-training} indicator of future PINN error, with a possible later
application to training. For residual
$r(t)=\dot{\widehat{x}}(t)-f(\widehat{x}(t))$, the hypothesis was that an
equally large residual should cause more later error if it points along a
locally expanding direction. The project tested alignment with the backward
Lyapunov vector $b_1$ and the leading finite-time singular vector $s_1$.
The observable prediction was a gain in held-out future-error prediction
beyond $\|r(t)\|$. The analysis plan fixed the horizons, controls and
decision rule before the original main test.

The submitted project report used Turinici's \emph{Regime-Aware Time
Weighting for Physics-Informed Neural Networks} as a reference. That paper
changes \emph{when} equation violations are weighted during training;
this project's main test asks whether residual \emph{orientation} is
informative after training. The two questions are related but distinct.

\section*{What was tested, and why the hypothesis was not supported}
The original experiment used Lorenz-63, a high-accuracy DOP853 reference,
analytic Jacobians and tangent dynamics, and 40 standard PINNs (eight
initial-condition windows, five seeds). It predicted future error on an
entirely held-out initial-condition window. Random, time-shuffled and
other dynamical directions were controls. Neither $b_1$ nor $s_1$ passed
the prespecified four-part rule at any usable horizon. Some within-run
associations failed to transfer to held-out windows or to clearly beat
controls. The proper conclusion is \emph{not supported for these tested
pointwise predictors}, not that direction can never matter.

Three observations help explain the result without proving a unique cause.
First, residuals were already highly aligned with $s_1$ (median alignment
about 0.96 at horizon 0.25), leaving little orientation variation. Second,
total future error inherits error from residuals produced earlier, which a
single current residual cannot summarize. Third, the original global PINNs
lost trajectory accuracy beyond a one-unit window, restricting the test.
An exploratory interpolation and trailing-average branch suggested that
recent residual history deserves study, but it did not change the original
verdict.

\section*{What changed in the model}
Exact initial-state enforcement improved the two-unit pilot but left median
trajectory error near six. Extra early-residual weighting and L-BFGS
refinement gave little additional improvement. The effective construction
was a \emph{sequential time-domain decomposition PINN}: eight separate
networks solve successive 0.25-unit intervals, each starting from the
previous network's \emph{predicted} endpoint. No true interior reference
state is used in training. In the development set, 35/40 runs stayed below
error one throughout two time units. The frozen model was then tested on a
new attractor run, with eight fresh evaluation windows and five seeds each.
The new reference and tangent computations passed all numerical checks.
On this independent set, 30/40 stayed below error one throughout and 32/40
for at least 90\% of the interval. The median of run-level median errors
was 0.0551. All five seeds on window 1 failed the full-window criterion,
and five further failures occurred across windows 3, 6 and 7. The model
is substantially better than the single global two-unit pilot, but less
reliable than the development-set rate suggested.

\section*{Current scientific status}
The original directional analysis was repeated on the 40 \emph{new} runs,
excluding exact segment joins where derivatives are two-sided. Neither
$b_1$ nor $s_1$ passed the original rule at any horizon. The same pattern
held when the analysis was restricted to accurate trajectory portions.
The oracle quantity
$\|\Phi(t+\Delta,t)r(t)\|$, which propagates the full residual vector
along the \emph{known reference trajectory}, increased held-out $R^2$ by
approximately $0.066$, $0.267$ and $0.369$ at horizons $0.25$, $0.5$ and
$1.0$. With elapsed time and its square in the baseline, whole-window
bootstrap gains were $+0.056$ [0.031, 0.083], $+0.258$ [0.206, 0.395]
and $+0.442$ [0.267, 0.672]. The oracle cannot be used as a deployable
predictor because it needs the true reference trajectory.

\begin{center}
\begin{tabular}{p{0.29\linewidth}p{0.62\linewidth}}
\toprule
Question & Current answer\\
\midrule
Original pointwise hypothesis & Not supported in the planned standard-PINN
test; the frozen follow-up repeated the pattern on fresh initial conditions.\\
Better Lorenz PINN & Two-unit accuracy below error one throughout in 30/40
new runs; accuracy depends on the initial condition.\\
Practical new error indicator & Not yet established. Full-propagator result is
oracle-only.\\
\bottomrule
\end{tabular}
\end{center}

\section*{When can the project reach a conclusion?}
\textbf{A project conclusion can be presented tomorrow:} the original
pointwise one-vector hypothesis was not supported for the tested Lorenz
PINNs, and a frozen repeat on fresh initial conditions showed the same
result. The sequential model extends two-unit accuracy to most, but not all,
new initial conditions. This is a bounded Lorenz-specific conclusion, not
proof that residual direction never matters. The sequential model was
designed after the first result, so its validation does not retroactively
become the original preregistration.

\textbf{A positive practical predictor remains open.} It would require a
causal statistic built from past PINN residuals alone, tested against
residual magnitude, time and simple trailing-history baselines. The
piecewise Duhamel quadrature also needs correction at segment joins.
Generality beyond Lorenz-63 would require another dynamical system.

\noindent\textbf{Evidence.} The frozen validation plan, chronological run
log and new 40-run accuracy figure are in the repository's \texttt{docs/}
directory. The original analysis plan and complete result arrays remain
available in the same project.

\end{document}
